% Folder 147, batch 10 — archivists' pages 181–189 (the last pages of the folder).
% Pass: Opus 5.5 (claude-opus-5-5), 2026-10-09 — first pass, unchecked against the pages by a human.
% All nine pages are mathematics in Grothendieck's hand; none skipped.
% His own pagination 87–90 appears on pp. 182, 184, 186, 188; p. 181 is the end of his p. 86.
\documentclass[11pt,a4paper]{article}
\input{../preamble/grothendieck}

\folder{147}
\batch{10}
\pages{181}{189}
\foldertitle{Structure à l'infini des Mg,ν [pages 1 à 90, dont table des matières] : notes manuscrites (s.d.).}
\dating{s.d. — le groupe « Autour de La "Longue Marche" à travers la théorie de Galois » (141 à 149) est daté [à partir de 1978]-1983}
\watermark{Édition de démonstration}

\begin{document}

\page{181}
\note{La page s'ouvre au milieu d'une phrase commencée avant ce lot (p.~86 de l'auteur) : il y est question de la structure locale d'une stratification au voisinage d'un point de la strate $X_i$.}

\marginal{NB les $Z_j$ loc. connexes}

isomorphe à une structure produit, $X_i \times Z$ (avec les $X_j$ correspondant à des $Z_j \subset Z$, $\overline{X_j}$ à $\overline{Z_j}$, $X_i$ à un $z_0 \in Z$ contenu dans les $\overline{Z_j}$), et les $\pi_0(Z_j\,;\overline{Z_j},z_0)$ sont des \add{\ill{}} cardinal $1$, i.e.\ il existe un syst. fond. de voisinages de $z_0$ dans $Z$ dont l'intersection avec les $Z_j$ est connexe)

\emph{Remarques} : Si $U$ est un ouvert de $X$, alors l'ens. des composantes connexes des $U \cap X_i$ forme \add{\ill{}} une stratification \add{globale} de $U$. Donc \struck{l'ens. l'app} \add{à un} ouvert $U$ variable de $X$,
\[
  U \longmapsto \mathrm{Strglob}(U)
\]
forme un préfaisceau sur $X$. Les exemples des d) montrent que ce préfaisceau n'est pas \struck{\ill{}} \add{un faisceau} (autre exemple encore plus simple : les feuilletages !)

\marginal{j'ai l'impression cependant qu'\ill{} stratification globale est connue quand on connaît \uncertain{localement}}

Les sections du faisceau associé s'appellent les \add{structures de} \emph{stratifications locales} sur $X$. Comparer avec l'exercice de SGA (théorie des Topos) sur les relations d'équivalence locales --- les stratifications en sont des exemples très

\page{182}
\note{p.~87 de l'auteur.}

particulières.

\struck{Si j'avais}

On suppose donnée sur $X$ une structure de stratification locale. \uncertain{Si j'en avais}, on pourrait comme dans le cas des feuilletages, plus gén. des toutes les relations d'équivalence locales) définir \struck{l'espace \add{déployé} des feuilles, que je vais} une partition de $X$ en des \struck{feuilles} \add{strates}, disons que \struck{feuilles} \add{Str. $E$} \uncertain{étant} munie d'une topologie (qui n'est pas en général la topologie induite) telle que l'application
\[
  E \longrightarrow X \tag{296}
\]
soit une immersion locale. Je vais désigner par $D^{*}_{\Delta_0}$ \struck{la \ill{}} \add{l'\uncertain{espace}} \ill{} des feuilles, donc
\[
  D^{*}_{\Delta_0} \longrightarrow X , \tag{297}
\]
\struck{\ill{}} est \uncertain{bijective}, et une imm. locale.

\note{Dans (296) le $E$ est écrit à double trait, sur un $D$ biffé ; dans (297) le $D$ est à double trait, sur un $E$ biffé. Plus loin il écrit $D$ tantôt simple, tantôt à double trait, sans différence apparente : on transcrit partout $D$. L'astérisque en exposant est souvent tracé comme un $\alpha$ (pp.~183, 185) ; le contexte impose $D^{*}$.}

\page{183}
\emph{Remarques} 1) Si $X$ est compact, on voit qu'il n'y a qu'un nb fini de \struck{feuilles} \add{strates} (qui sont loc. compactes) --- on s'intéresse surtout aux \struck{feuill\supplied{etages}} \add{stratifications} avec un nb fini de strates --- mais bien sûr cette hyp. n'est pas stable par passage à un ouvert de $X$.

b) La formation de $D^{*}_{\Delta_0}$ commute à la restriction à un ouvert de $X$.

c) Si la \struck{\ill{}} \add{structure} de strat. est définie par un diviseur à croisements normaux dans un \uncertain{espace} analytique, alors on a un homéomorphisme
\[
  D^{*}_{\Delta_0} \simeq \coprod_{d \in \mathbb{N}} D^{*}_{d} \tag{298}
\]
donc la notation $D^{*}_{\Delta_0}$ introduite ici coïncide avec celle des \struck{\ill{}} deux sections précédentes.

\page{184}
\note{p.~88 de l'auteur ; en tête, à côté du numéro, un mot biffé illisible.}

On appelle \emph{drapeau} de longueur $r \in \mathbb{N}$ d'une stratification globale tout système
\[
  \Sigma_* = \{\overline{X_{i_0}} \subset \overline{X_{i_1}} \subset \dots \subset \overline{X_{i_r}}\} \tag{299}
\]
de strates --- le drapeau est \emph{non dégénéré} si les inclusions sont strictes. Un drapeau \emph{local} \add{\emph{strict}} de longueur $r$ est un \struck{système} couple formé d'un tel drapeau, et d'une « origine » \struck{$x \in X_{i_0}$ \emph{NB}}.
\[
  x \in X_{i_0} \tag{300}
\]
\emph{NB} on exige $x \in X_{i_0}$ et non seulement $x \in \overline{X_{i_0}}$.

\struck{On met une topologie sur l'ens. des drapeaux du type donné, en \ill{} les drapeaux} \note{passage encadré et barré en croix.}

On met une topologie sur l'ens. des drapeaux \add{locaux stricts} \struck{(de type)}, de longueur $r$, en prenant comme \add{base d'ouverts les} \struck{voisinages des $(x,\Sigma_*)$} ensembles qu'on obtient en prenant $\Sigma_*$ fixe, et $x$ variable dans un \uncertain{ouvert} de $\Sigma_*$. On trouve un espace topologique $D^{*}_{\Delta_r}$, qui s'envoie dans $X$ par

\page{185}
\marginal{les images des composantes connexes des $D^{*}_{\Delta_r}$ sont des strates}

$(x,\Sigma_*) \longmapsto x$, d'où
\[
  D^{*}_{\Delta_r} \longrightarrow X \tag{301}
\]
qui est une immersion locale. En fait, cet espace sur $X$ ne dépend que de la \add{structure de} stratification locale sur $X$ définie par la stratification globale --- on pourrait la définir directement en termes de la structure locale à l'aide de la notion des « drapeaux en germes de strates ». \struck{Une autre façon de procéder serait de} \struck{On aura} \add{Dans} le cas d'une stratification associée à un diviseur à cr. norm. sur un esp. analytique, on retrouve les notations déjà utilisées
\[
  D^{*}_{\Delta_r} = \coprod_{d_r \in \mathcal{P}_{r+1}(\mathbb{N})} D^{*}_{d_r}
\]
des $D_{\Delta_r}$ des deux sections précédentes. Mais il faudrait pouvoir définir l'analogue des morphismes ½simpliciaux

\note{Le $\Delta_r$ de gauche est écrit sur un symbole biffé ; sous le $\coprod$, un indice biffé.}

\page{186}
\note{p.~89 de l'auteur.}
\[
  D^{*}_{\Delta_r} \longrightarrow D_{\Delta_{r'}} \quad \text{pour } \Delta_{r'} \to \Delta_r
\]
et pour ceci, on \struck{doit} \add{pourrait} définir quelque chose qui tienne lieu de $D_{\Delta_{r'}}$ ---

On va définir les drapeaux \add{(pas stricts)} \struck{généraux} comme étant les systèmes $(x,\Sigma_*)$, mais cette fois-ci \add{seulement}
\[
  x \in \overline{X_{i_0}} .
\]
On définit la topologie comme précédemment, on trouve un $D_{\Delta_r}$ et
\[
  D_{\Delta_r} \longrightarrow X , \tag{302}
\]
qui est encore une immersion locale, et $D^{*}_{\Delta_r}$ est un ouvert canonique de $D_{\Delta_r}$.

Cette fois-ci on a fait ce qu'il fallait pour que les $D_{\Delta_r}$ forment un système ½simplicial d'espaces --- les morphismes de dégénérescence fournissant des \ill{}.

Cet espace ½simplicial est-il

\page{187}
une catégorie dans $(\mathrm{Top})$ ? Oui, en d'autres termes, l'application canonique
\[
  D_{\Delta_2} \xrightarrow{\ \sim\ } D_{\Delta_1} \times_{D_{\Delta_0}} D_{\Delta_1} \tag{303}
\]
est iso --- cela exprime simplement que la donnée d'un drapeau \add{local} de longueur $2$
\[
  \Sigma_* = \{x \in \overline{X_{i_0}} \subset \overline{X_{i_1}} \subset \overline{X_{i_2}}\}
\]
revient à la donnée de $\{x \in \overline{X_{i_0}} \subset \overline{X_{i_1}}\}$ et $\{x \in \overline{X'_{j_0}} \subset \overline{X'_{j_1}}\}$ telles que $\overline{X_{i_1}} = \overline{X'_{j_0}}$ (il s'agit de germes autour de $x$…)

Donc une \emph{stratification locale} définit une « catégorie topologique »
\[
  \left\{
  \begin{aligned}
    &\text{espace des objets} = D_{\Delta_0} \\
    &\qquad = \text{espace des strates locales pas strictes} \\
    &\text{espace des flèches} = D_{\Delta_1} \\
    &\qquad = \text{espace des drapeaux } x \in \overline{X_{i_0}} \subset \overline{X_{i_1}}
  \end{aligned}
  \right. \tag{303}
\]
\struck{morphismes} « source », « but », « composition » définis de façon évidente.

\note{Le numéro (303) est porté deux fois sur cette page, par l'auteur. Devant « pas strictes », un mot biffé illisible.}

\page{188}
\note{p.~90 de l'auteur.}

On a ceci :
\[
  D_{\Delta_r} \xrightarrow{\ \sigma\ } D_{\Delta_0} \tag{304}
\]
\marginal{application simpliciale correspondant au « premier sommet » d'un simplexe ordonné ; ou à la « source » d'une suite de flèches $\bullet \to \bullet \to \bullet \cdots \to \bullet$.}

\marginal{\emph{NB} $D^{*}_{\Delta_r}$ est l'image inverse de $D^{*}_{\Delta_0}$}

\struck{est} une application finie étale \emph{au dessus} de $D^{*}_{\Delta_0}$ (\emph{NB} Dans le cas « \ill{} \ill{} », c'est \ill{} vrai \ill{} cette fois sur $D_{\Delta_0}$ tout entier --- mais c'est là un fait très spécial, dû au fait que les \struck{\ill{}} strates sont des variétés topologiques --- ce qui implique les \ill{} pour \ill{} \ill{} des $D_{\Delta_\bullet}$…)

Également le fait que \struck{\ill{}} la structure \add{locale} ordinaire \add{de l'une} des $D^{*}_{\Delta_r}$ sur $D^{*}_{\Delta_0}$ \struck{(flèches} structures \add{des} strates locales contenues dans une strate locale donnée) soit \struck{\ill{}} \add{\ill{}} \struck{structure de simplexe}

\page{189}
à l'une des parties d'un ens. fini, est spécial aux croisements normaux et n'est déjà plus le cas pour
\drawing{un point d'où partent trois arcs (un « tripode »)}
où le nb des strates contenant $\bullet$ \ill{} \ill{} $=$ \struck{\ill{}} $7$, qui n'est pas une puissance de deux…

\emph{NB} Pour l'instant, on n'a pas \add{\uncertain{du tout}} fait usage de l'hypothèse \add{de local trivialité} de p.~86 e). Elle servira plus loin, pour expliciter en termes de fibrations loc. triviales certains voisinages tubulaires, \struck{\ill{}} \add{\uncertain{fournissant} \ill{}} des « types élémentaires » de la section précédente.

\note{Fin du dossier : la p.~90 de l'auteur est la dernière ; l'argument reste ouvert.}

\end{document}
